% ------------------------------------------------------------
% Preamble for 「Pythonによる はじめての銀河シミュレーション」
% 157 x 221 mm layout based on the supplied textbook template.
% Compile with XeLaTeX.
% ------------------------------------------------------------
\usepackage[paperwidth=157mm,paperheight=221mm,margin=14mm,headheight=14pt,headsep=6mm]{geometry}
\usepackage{amsmath,amssymb,mathtools,bm}
\usepackage{graphicx,rotating,float,placeins,caption}
\captionsetup{font=small,labelfont=bf,labelsep=quad}
\graphicspath{{figures/}{./}}
\usepackage{fontspec}
\usepackage{xeCJK}
\IfFontExistsTF{Latin Modern Roman}{\setmainfont{Latin Modern Roman}}{\setmainfont{Liberation Serif}}
\IfFontExistsTF{Latin Modern Sans}{\setsansfont{Latin Modern Sans}}{\setsansfont{Liberation Sans}}
\IfFontExistsTF{DejaVu Sans Mono}{\setmonofont{DejaVu Sans Mono}[Scale=0.91]}{}
\setCJKmainfont{Noto Serif CJK JP}[ItalicFont={Noto Serif CJK JP}]
\setCJKsansfont{Noto Sans CJK JP}
\setCJKmonofont{Noto Sans Mono CJK JP}[Scale=0.90]
\xeCJKsetup{CJKmath=true}
\XeTeXlinebreaklocale "ja"
\XeTeXlinebreakskip = 0pt plus 0.1pt minus 0.01pt

\usepackage{xcolor,enumitem,array,booktabs,tabularx,xurl,etoolbox,needspace,fvextra}
\usepackage{fix-cm}

% TikZ settings are defined here rather than in a separate input file.
% This makes the project robust on Overleaf even when only the main preamble
% and chapter/figure files are copied into an existing project.
\usepackage{tikz}
\usetikzlibrary{arrows.meta,calc,positioning,patterns,decorations.pathreplacing,matrix}
\definecolor{figblue}{HTML}{0072B2}
\definecolor{figred}{HTML}{D55E00}
\definecolor{figgreen}{HTML}{009E73}
\definecolor{figgray}{HTML}{666666}
\tikzset{
  textbook diagram/.style={font=\sffamily\fontsize{9}{11}\selectfont,>=Latex,line width=0.6pt},
  vector/.style={-{Latex[length=2.1mm,width=1.4mm]},line width=0.85pt},
  guide/.style={draw=figgray,densely dashed,line width=0.5pt},
  state/.style={draw=black!50,fill=white,rounded corners=1pt,minimum width=48mm,minimum height=9mm,align=center},
  hint/.style={align=left,font=\sffamily\fontsize{8.5}{10.5}\selectfont},
  small label/.style={font=\sffamily\fontsize{8}{10}\selectfont}
}

\usepackage[most]{tcolorbox}
\usepackage{listings}
\usepackage{titlesec,fancyhdr}

\setlength{\parindent}{1em}
\setlength{\parskip}{0.15em}
\linespread{1.08}
\raggedbottom
\emergencystretch=1.5em
\widowpenalty=10000
\clubpenalty=10000

\definecolor{navy}{HTML}{0D1B3E}
\definecolor{amber}{HTML}{B5740F}
\definecolor{amberbg}{HTML}{FBF1E0}
\definecolor{cyan}{HTML}{1F6390}
\definecolor{cyanbg}{HTML}{E7F0F7}
\definecolor{navybg}{HTML}{EDF0F7}
\definecolor{rule}{HTML}{C9D2E3}
\definecolor{notegray}{HTML}{F3F5F6}
\definecolor{notegreen}{HTML}{287250}
\definecolor{warnred}{HTML}{9B4030}

% Article boxes: notebox/failbox. goalbox is a chapter-opening structural element.
\newtcolorbox{notebox}[1]{enhanced,breakable,colback=notegray,colframe=notegreen,
  boxrule=0.45pt,arc=1.5pt,left=2.8mm,right=2.8mm,top=2mm,bottom=2mm,
  fonttitle=\sffamily\bfseries,colbacktitle=notegreen!10,coltitle=notegreen,
  title={もう少し詳しく：#1},before skip=7pt,after skip=7pt}
\newtcolorbox{failbox}[1]{enhanced,breakable,colback=notegray,colframe=warnred,
  boxrule=0.45pt,arc=1.5pt,left=2.8mm,right=2.8mm,top=2mm,bottom=2mm,
  fonttitle=\sffamily\bfseries,colbacktitle=warnred!8,coltitle=warnred,
  title={よくある失敗：#1},before skip=7pt,after skip=7pt}
\newtcolorbox{goalbox}[1][本章の目標]{colback=cyanbg,colframe=cyan,coltitle=white,
  title={\sffamily #1},fonttitle=\bfseries,boxrule=0.7pt,arc=2pt,
  left=7pt,right=7pt,top=5pt,bottom=5pt,breakable}

\newcommand{\inputcode}[1]{\par\smallskip\noindent\begin{minipage}{\linewidth}\VerbatimInput[
  fontsize=\fontsize{7.5}{9.3}\selectfont,frame=lines,rulecolor=\color{navy!35},framesep=2.2mm,
  numbers=left,numbersep=4pt,breaklines=true,breakanywhere=true,tabsize=4]{generated/#1.py}\end{minipage}\par\smallskip}

\titleformat{\chapter}{\LARGE\bfseries\sffamily\color{navy}}{第\thechapter 章}{0.8em}{}
\titleformat{\section}{\large\bfseries\sffamily\color{navy}}{\thesection}{0.5em}{}
\titleformat{\subsection}{\normalsize\bfseries\sffamily\color{navy}}{\thesubsection}{0.5em}{}
\titlespacing*{\chapter}{0pt}{0pt}{12pt}
\titlespacing*{\section}{0pt}{1.1em}{0.4em}
\titlespacing*{\subsection}{0pt}{0.75em}{0.25em}
\pretocmd{\section}{\Needspace{6\baselineskip}}{}{}
\pretocmd{\subsection}{\Needspace{5\baselineskip}}{}{}
\setlist[itemize]{leftmargin=1.7em,itemsep=0.15em,topsep=0.25em}
\setlist[enumerate]{leftmargin=2em,itemsep=0.2em,topsep=0.3em}
\setcounter{tocdepth}{1}
\setcounter{secnumdepth}{2}

\pagestyle{fancy}
\fancyhf{}
\fancyhead[L]{\scriptsize\sffamily はじめての銀河シミュレーション}
\fancyhead[R]{\scriptsize\sffamily\nouppercase{\leftmark}}
\makeatletter
\renewcommand{\chaptermark}[1]{%
  \if@mainmatter
    \markboth{第\thechapter 章\enspace #1}{}%
  \else
    \markboth{#1}{}%
  \fi}
\makeatother
\fancyfoot[R]{\small\thepage}
\fancypagestyle{plain}{\fancyhf{}\fancyfoot[R]{\small\thepage}\renewcommand{\headrulewidth}{0pt}}

\newcommand{\rvec}{\boldsymbol r}
\newcommand{\vvec}{\boldsymbol v}
\newcommand{\avec}{\boldsymbol a}
\newcommand{\yvec}{\boldsymbol y}
\newcommand{\fvec}{\boldsymbol f}
\newcommand{\kvec}{\boldsymbol k}
\newcommand{\Rvec}{\boldsymbol R}
\newcommand{\Vvec}{\boldsymbol V}
\newcommand{\Dvec}{\boldsymbol D}
\newcommand{\Uvec}{\boldsymbol U}
\newcommand{\Pvec}{\boldsymbol P}
\newcommand{\Mtot}{M_{\rm tot}}
\newcommand{\mred}{m_{\rm red}}
\newcommand{\Tc}{T_{\rm c}}
\newcommand{\Msun}{M_\odot}
\newcommand{\unit}[1]{\,\mathrm{#1}}
\newcommand{\dd}{\mathrm{d}}
\newcommand{\vect}[1]{\boldsymbol{#1}}
\newcommand{\figpath}[1]{figures/#1.pdf}
\renewcommand{\contentsname}{目次}
\renewcommand{\figurename}{図}
\renewcommand{\tablename}{表}
\renewcommand{\bibname}{参考図書}
% Chapter zero must keep the prefix in figures, tables, and equations.
\renewcommand{\thefigure}{\thechapter.\arabic{figure}}
\renewcommand{\thetable}{\thechapter.\arabic{table}}
\renewcommand{\theequation}{\thechapter.\arabic{equation}}
% Same end-of-chapter headings and bibliography in every chapter.
\newcommand{\chapterexercises}{\section*{演習}\phantomsection\addcontentsline{toc}{section}{演習}}
\newcommand{\answerchapter}[2]{\section*{第#1章　#2}\phantomsection\addcontentsline{toc}{section}{第#1章　#2}}
% Shell commands are not executable Python; keep them visibly distinct.
\newcommand{\inputterminal}[1]{\par\smallskip\noindent
  {\footnotesize\sffamily 端末で実行}\par\nobreak
  \VerbatimInput[fontsize=\fontsize{8}{10}\selectfont,frame=single,rulecolor=\color{black!30},
  framesep=2mm,breaklines=true,breakanywhere=true,numbers=none]{#1}\smallskip}
\newcommand{\booktitle}{Pythonによる はじめての銀河シミュレーション}
\newcommand{\booksubtitle}{テスト粒子で学ぶ重力・軌道・銀河相互作用}
\newcommand{\bookversion}{見出し統一版 v2.5}
\newcommand{\distributionnote}{授業・ゼミで案内する配布先から入手してください。公開URLは刊行時に確定します。}


\renewcommand{\topfraction}{0.95}
\renewcommand{\bottomfraction}{0.85}
\renewcommand{\textfraction}{0.05}
\renewcommand{\floatpagefraction}{0.85}
\setcounter{topnumber}{3}
\setcounter{totalnumber}{5}
\setlength{\textfloatsep}{9pt plus 3pt minus 2pt}
\setlength{\intextsep}{9pt plus 3pt minus 2pt}

\usepackage[unicode,colorlinks=true,linkcolor=cyan,citecolor=cyan,urlcolor=cyan]{hyperref}
